% =====================================================================
% Capítulo 5 — Conclusiones
% =====================================================================
\chapter{Conclusiones}
\label{cap:conclusiones}

\section{Resumen del trabajo realizado}
\label{sec:resumen}

\todo{Recapitular brevemente: revisión de fundamentos teóricos
de los modelos de difusión, experimento controlado con mariposas,
diseño e implementación del sistema de reconstrucción PET en dos
\textit{pipelines} (supervisado condicional y prior no
condicional con DPS) más dos líneas base discriminativas (U-Net
de regresión y RED-CNN), infraestructura reproducible en la nube,
evaluación comparativa sobre el conjunto de test y análisis de
costes.}

\section{Resultados principales}
\label{sec:resultados}

\todo{Sintetizar los resultados cuantitativos (PSNR, SSIM,
NRMSE, preservación de intensidad) y cualitativos (figuras
4-panel) obtenidos por ambos \textit{pipelines}, y discutir la
comparación entre el enfoque supervisado y el prior no
condicional con guiado.}

\todo{Mostrar los volúmenes NIfTI reconstruídos completamente
por el pipeline de difusión, y compararlos con los volúmenes 3d originales}

\section{Limitaciones}
\label{sec:limitaciones}

La mayor limitación del proyecto es la combinación entre la
demanda de máquinas de alto rendimiento para la experimentación y
la amplitud del espacio de experimentos posibles. Este espacio
abarca desde los parámetros del paradigma de entrenamiento
(tamaño de imagen, tasa de aprendizaje, estabilización numérica),
pasando por las decisiones de tratamiento de las imágenes médicas
(normalización, recorte, filtrado de \textit{slices}), hasta las
elecciones de arquitectura (DDIM frente a DDPM, U-Net con
atención, guiado DPS, imitación explícita del ruido de dosis
baja, etc.). Combinado con los límites de hardware, que solo
permiten realizar un número acotado de experimentos, esto ha sido
el mayor reto del trabajo. Se ha mitigado mediante una revisión
del estado del arte exhaustiva, de la que se han extraído las
ideas más prometedoras de la literatura reciente del dominio; sin
ella habría sido muy difícil llegar a tiempo, pues habría hecho
falta una experimentación mucho mayor para determinar qué funciona,
qué no, y dónde se rompe la arquitectura.

Una segunda limitación es el tiempo de reconstrucción de un
volumen NIfTI en inferencia. Tras entrenar el modelo supervisado
hasta el final (100 épocas) y reconstruir volúmenes localmente en
un MacBook Pro M4 Pro de 24~GB con un \textit{batch} de inferencia
de 8, el tiempo de reconstrucción de un volumen completo es de
aproximadamente una hora. Esto dificulta integrar el sistema
directamente en una máquina de escaneo, tanto por los
requerimientos computacionales como por el tiempo de procesado.
De cara a un entorno clínico real, una arquitectura factible sería
adquirir el escaneo de dosis baja en la propia máquina y enviar
los volúmenes en lotes a la nube, donde se procesarían de forma
concurrente sobre sistemas de alto rendimiento. Esta solución
mantendría intacta la premisa del proyecto: obtener escaneos
equivalentes a los de dosis completa de radiación empleando una
vigésima parte de la dosis.

\section{Líneas de desarrollo futuro}
\label{sec:lineas-futuras}

Se identifican ocho extensiones, ordenadas por proximidad a la
implementación actual.

\begin{enumerate}[leftmargin=*]
    \item \textbf{Mayor resolución de entrenamiento ($384^2$).}
    Sustituir los \textit{slices} de $256^2$ empleados en el
    entrenamiento por \textit{slices} de $384^2$, una resolución
    más próxima a los $440^2$ nativos del volumen
    (\Cref{sec:pet-preprocesado}). Al preservar más detalle
    espacial cabe esperar un modelo más preciso, cuyas
    reconstrucciones se acerquen más a la referencia de dosis
    completa, a costa de un mayor consumo de memoria y cómputo por
    \textit{slice}. Es la extensión más inmediata, pues se reduce
    a un cambio de la resolución objetivo en la configuración.

    \item \textbf{Refinamiento de las líneas base U-Net y RED-CNN.}
    Alterar la arquitectura de los dos modelos discriminativos de
    referencia y, sobre todo, sustituir su objetivo de
    entrenamiento por una pérdida más rica que combine varias
    métricas complementarias (p.\,ej.\ una mezcla ponderada de MAE,
    MSE y un término estructural tipo SSIM), en lugar de un único
    error puntual. Dado que estas líneas base ya resultan
    ligeramente superiores a los \textit{pipelines} de difusión, un
    objetivo multi-métrica podría ampliar esa ventaja y precisar
    mejor el techo alcanzable por los métodos discriminativos.

    \item \textbf{DPS con $\sigma(\mathbf{y})$ dependiente de la
    intensidad.} Modelar la estadística de Poisson del ruido PET
    propagada a través de la normalización $\operatorname{arcsinh}$
    y sustituir la verosimilitud isotrópica del Pipeline~B por
    una verosimilitud heteroscedástica. Es la extensión natural
    más cercana al planteamiento sobre el ruido descrito en la
    \Cref{sec:pet-ruido}.

    \item \textbf{Replanteamiento del \textit{pipeline} no
    condicional mediante modelado explícito del ruido del
    escáner.} El guiado a posteriori (DPS) no ha producido
    resultados eficaces y el uso de ruido gaussiano del
    \textit{schedule} no captura la naturaleza del ruido PET. Se
    propone abandonar esa estrategia en favor de modelar
    explícitamente el ruido del escáner con técnicas alternativas
    —desde verosimilitudes ajustadas a la estadística de Poisson
    hasta procesos directos que imiten la degradación de dosis
    baja (véase el ítem sobre modelos puente)—, de modo que el
    proceso generativo aprenda la transición real entre dosis
    completa y dosis baja en lugar de una corrupción gaussiana
    sintética.

    \item \textbf{Contexto 2.5D o 3D.} Añadir \textit{slices}
    vecinos como canales adicionales de entrada (2.5D) para
    introducir coherencia inter-\textit{slice}, o pasar a un
    U-Net 3D sobre parches si el hardware se amplía.

    \item \textbf{Reponderado de pérdida por SNR mínimo.}
    Activación del \textit{Min-SNR-$\gamma$} de Hang et
    al.~\cite{hang2023minsnr} si la convergencia del
    \textit{v-prediction} con 100 épocas resulta insuficiente.

    \item \textbf{Modelos puente (I\textsuperscript{2}SB).}
    Liu et al.~\cite{liu2023i2sb} proponen entrenar un modelo
    cuyo proceso directo mapee directamente la distribución de
    dosis completa a la de dosis baja mediante un puente de
    Schrödinger, sustituyendo el ruido gaussiano del
    \textit{schedule} por una transición entre las dos
    distribuciones de interés. Es la línea más cercana a la
    observación recogida en la \Cref{sec:pet-ruido}: el proceso
    directo del modelo imita explícitamente la reducción de
    dosis. Constituye una línea prometedora para un capítulo
    de \textit{contribución} en la memoria, o como punto de
    partida para una futura tesis doctoral.

    \item \textbf{Métricas clínicas.} En el caso de obtener
    metadatos de dosis inyectada y peso del paciente, computar
    SUV reales y reportar preservación de SUV en regiones de
    interés segmentadas (manualmente o con un modelo auxiliar).
\end{enumerate}

\section{Reflexión final}
\label{sec:reflexion}

\todo{Reflexión personal sobre el aprendizaje adquirido durante
el TFM y su encaje con los contenidos del Máster en Tratamiento
Estadístico-Computacional de la Información.}
